\problem{}

This problem is a nice roadmap to how a theory is canonically quantized. 
We start of from the classical theory defined by it's classical Lagrangian 
then compute all of it's quantities of interest, still classically, and 
lastly promote the fields to operators by imposing the canonical commutation 
relations, which also amounts to follow the prescription $\acomm{\cdot}{\cdot}\rightarrow \frac{1}{\im\hbar}\comm{\cdot}{\cdot}$. 
Of course, here everything ``works out'' nicely, as the theory under consideration 
is a free theory, those are definable non-perturbatively. When we get to 
interacting theories not everything will function out of the box, but, 
many of the developments here will be useful.

\probitem{}

For the canonical quantization prescription, the first step is to go to 
the Hamiltonian phase-space formalism, in which lives the Poisson bracket, 
compute the conjugated momenta, and reexpress all the quantities of interest 
in terms of the fields and their momenta. As we know, the definition of the 
canonical conjugated momenta is simply the derivative of the lagrangian with 
respect to the time derivative of the fields,
\begin{align}
    \pi^I  = \pdv{\mathcal L}{\dot \phi_I},\ \ \ \dot\phi_ I = \partial_0\phi_ I
\end{align}
for our choice of Lagrangian,
\begin{align}
    \mathcal L = -\partial_\mu\phi^\ast\partial^\mu\phi-m^2\phi^\ast\phi
\end{align}
there are two kinds of dynamical objects appearing, $\phi,\phi^\ast$. It's clear that 
$\phi^\ast$ is the complex conjugate of $\phi$. In principle one could argue that we have just 
one dynamical field, $\phi$, and that $\phi^\ast$ is completely determined by $\phi$. 
But this reasoning is not true, as a complex field,  
each $\phi$ has two real degrees of freedom, nevertheless, is not possible to 
extract those directly from $\phi$ alone. For us to define the real and 
imaginary parts we need a definition of $\phi^\ast$,
\begin{align}
    \mathfrak{Re}\qty[\phi]=\frac12\qty(\phi+\phi^\ast),\ \ \ \mathfrak{Im}\qty[\phi] = -\im\frac12\qty(\phi-\phi^\ast)
\end{align}
also, from complex variables calculus, we also know that $z,z^\ast$ are 
independent quantities, which can be seen from relations of the kind,
\begin{align}
    \pdv{ z^\ast}{z} = 0
\end{align}
here is the same principle.

All of this is just to argue that, in the expressions where we included a 
$\phi_I$ to denote multiple fields, there are two possible choices of fields, 
$\phi_I\in\qty{\phi,\phi^\ast}$, that means that are two canonical conjugated 
momenta, one for each field\footnote{Actually the discussion is a lot more subtle to how the 
canonical momenta are defined, and which fields should have one associated with.},
\begin{align}
    \pi  = \pdv{\mathcal L}{\dot \phi},\ \ \ \pi^\ast  = \pdv{\mathcal L}{\dot \phi^\ast}
\end{align}
it's straightforward to obtain then from the Lagrangian,
\begin{align}
    \pi = \dot\phi^\ast,\ \ \ \pi^\ast = \dot\phi
\end{align}

The next step is to proceed with the computation of the Hamiltonian 
density, which is given by the Legendre transform of the Lagrangian,
\begin{align}
    \mathcal H  &= \pi\dot \phi+\pi^\ast\dot \phi^\ast - \mathcal L\\
    \mathcal H  &= \pi\pi^\ast+\pi^\ast\pi +\partial_\mu\phi^\ast\partial^\mu\phi +m^2\phi^\ast\phi\\
    \mathcal H  &= 2\pi^\ast\pi -\dot\phi^\ast\dot \phi+\grad\phi^\ast\cdot\grad\phi +m^2\phi^\ast\phi\\
    \mathcal H  &= 2\pi^\ast\pi -\pi\pi ^\ast+\grad\phi^\ast\cdot\grad\phi +m^2\phi^\ast\phi\\
    \mathcal H  &= \pi^\ast\pi+\grad\phi^\ast\cdot\grad\phi +m^2\phi^\ast\phi
\end{align}
and finally to compute the full Hamiltonian is only necessary to integrate in the spatial variables,
\begin{align}
    H&=\int\dd[3]{\vb x}\mathcal H\\
    H&=\int\dd[3]{\vb x}\qty[\pi^\ast\pi+\grad\phi^\ast\cdot\grad\phi +m^2\phi^\ast\phi]
\end{align}

\probitem{}

This is a nice toy model symmetry which generates a charge in the proper 
sense. As we'll see in due time, the electron ``charge'' arises due to 
the existence of a phase symmetry exactly as this one. As we commented 
before, by $\phi$ being a complex field, there are two degrees of freedom, 
this will imply in two particles in our theory, they'll have opposite 
sign of this charge generated by the phase symmetry. The same happens to 
the electron and positron.

First let's check it is a symmetry, the transformation law is,
\begin{align}
    \begin{cases}
        x&\rightarrow x' = x\\
        \phi(x)&\rightarrow \phi'\qty(x') = \e^{-\im a}\phi\qty(x)\\
        \phi^\ast(x)&\rightarrow {\phi^\ast}'\qty(x') = \e^{\im a}\phi^\ast\qty(x)
    \end{cases}\Rightarrow \begin{cases}
        \dv{\phi'(x)}{a}\eval_{a=0} = -\im\phi\\
        \dv{{\phi^\ast}'(x)}{a}\eval_{a=0} = \im\phi
    \end{cases}
\end{align}
substituting back in the action,
\begin{align}
    S_a'[\phi']&=\int\dd[4]{x'}\qty(-\partial_\mu '{\phi^\ast}'{\partial^\mu}'\phi'-m^2{\phi^\ast}'\phi')\\
    S_a'[\phi']&=\int\dd[4]{x}\qty(-\partial_\mu \e^{\im a}{\phi^\ast}{\partial^\mu}\e^{-\im a}\phi-m^2\e^{\im a}{\phi^\ast}\e^{-\im a}\phi)\\
    S_a'[\phi']&=\int\dd[4]{x}\qty(-\partial_\mu{\phi^\ast}{\partial^\mu}\phi-m^2{\phi^\ast}\phi) = S[\phi]
\end{align}
indeed is a symmetry. Then, we can use Noether's theorem,
\begin{align}
    \tensor{J}{^\mu} &=\pdv{\mathcal L}{\partial_\mu\phi_ I\qty(x)}\dv{\phi_ I'\qty(x)}{a}\eval_{a=0}+\mathcal L\dv{{x'}^\mu}{a}\eval_{a=0}
\end{align}
as we commented before, we have to take $\phi_I=\qty{\phi,\phi^\ast}$,
\begin{align}    
    \tensor{J}{^\mu} &=\pdv{\mathcal L}{\partial_\mu\phi\qty(x)}\dv{\phi'\qty(x)}{a}\eval_{a=0}+\pdv{\mathcal L}{\partial_\mu\phi^\ast\qty(x)}\dv{{\phi^\ast}'\qty(x)}{a}\eval_{a=0}\\
    \tensor{J}{^\mu} &=-\partial^\mu\phi^\ast\dv{\phi'\qty(x)}{a}\eval_{a=0}-\partial^\mu\phi\dv{{\phi^\ast}'\qty(x)}{a}\eval_{a=0}\\
    \tensor{J}{^\mu} &=-\partial^\mu\phi^\ast\qty(-\im\phi)-\partial^\mu\phi\qty(\im \phi^\ast)\\
    \tensor{J}{^\mu} &=-\im \phi^\ast\partial^\mu\phi+\im\phi\partial^\mu\phi^\ast
\end{align}
to obtain the conserved charge we integrate the temporal component over a spatial slice,
\begin{align}
    Q&= \int\dd[3]{\vb x}J^0=\int\dd[3]{\vb x}\qty(-\im \phi^\ast\partial^0\phi+\im\phi\partial^0\phi^\ast)\\
    Q&=\int\dd[3]{\vb x}\qty(\im \phi^\ast\pi^\ast-\im\phi\pi)\\
    Q&=-\im\int\dd[3]{\vb x}\qty(\pi\phi-\pi^\ast \phi^\ast)
\end{align}
as we'll see in the remaining of the problems, the minus sign between $\pi\phi$ and $\pi^\ast\phi^\ast$ is 
what gives a difference of sign in the charges of particles and anti-particles of 
this theory.

\probitem{}

This serves as a starting point of the quantization, showing the quantum equations of motion, of, the Heisenberg equations. 
As we know, these are the analogous classical equations of motion given by the Poisson bracket in the Hamiltonian formalism,
\begin{align}
    \acomm{A\qty(t,\vb x)}{H\qty(t)}=\dot A\qty(t,\vb x)\Rightarrow \comm{A\qty(t,\vb x)}{ H(t)} = \im\dot A\qty(t,\vb x) 
\end{align}

We already found an expression for the Hamiltonian. What is missing is the canonical commutation relations. These also mimic the 
canonical Poisson brackets, let's state them here, 
\begin{align}
    \comm{\phi\qty(t,\vb x)}{\pi\qty(t,\vb y)}&= \im\delta^{\qty(3)}\qty(\vb x-\vb y), \ \ \ \comm{\phi\qty(t,\vb x)}{\phi\qty(t,\vb y)}=\comm{\pi\qty(t,\vb x)}{\pi\qty(t,\vb y)}=0\\
    \comm{\phi^\dagger\qty(t,\vb x)}{\pi^\dagger\qty(t,\vb y)}&= \im\delta^{\qty(3)}\qty(\vb x-\vb y), \ \ \ \comm{\phi^\dagger\qty(t,\vb x)}{\phi^\dagger\qty(t,\vb y)}=\comm{\pi^\dagger\qty(t,\vb x)}{\pi^\dagger\qty(t,\vb y)}=0\\
    \comm{\phi^\dagger\qty(t,\vb x)}{\phi\qty(t,\vb y)}&=\comm{\pi^\dagger\qty(t,\vb x)}{\pi\qty(t,\vb y)}=\comm{\phi^\dagger\qty(t,\vb x)}{\pi\qty(t,\vb y)}=0
\end{align}
which we can use to compute the commutator of $H$ and $\phi$, with the already found 
expression of $H$,
\begin{align}
    \im\dot\phi\qty(t,\vb x)&=\comm{\phi\qty(t,\vb x)}{H\qty(t)}\\
    \im\dot\phi\qty(t,\vb x)&=\comm{\phi\qty(t,\vb x)}{\int\dd[3]{\vb y}\qty[\pi^\dagger\qty(t,\vb y)\pi\qty(t,\vb y)+\grad\phi^\dagger\qty(t,\vb y)\cdot\grad\phi\qty(t,\vb y) +m^2\phi^\dagger\qty(t,\vb y)\phi\qty(t,\vb y)]}\\
    \im\dot\phi\qty(t,\vb x)&=\int\dd[3]{\vb y}\comm{\phi\qty(t,\vb x)}{\pi^\dagger\qty(t,\vb y)\pi\qty(t,\vb y)+\grad\phi^\dagger\qty(t,\vb y)\cdot\grad\phi\qty(t,\vb y) +m^2\phi^\dagger\qty(t,\vb y)\phi\qty(t,\vb y)} 
\end{align}
from the three terms inside the commutator, only the first one will give a non-zero result, 
due to the canonical commutation relations,
\begin{align}
    \comm{\phi\qty(t,\vb x)}{\phi\qty(t,\vb y)} = \comm{\phi\qty(t,\vb x)}{\phi^\dagger\qty(t,\vb y)} =0\Rightarrow \comm{\phi\qty(t,\vb x)}{\grad\phi\qty(t,\vb y)} = \comm{\phi\qty(t,\vb x)}{\grad\phi^\dagger\qty(t,\vb y)} =0
\end{align}
with this is straightforward to finish the calculation,
\begin{align}
    \im\dot\phi\qty(t,\vb x)&=\int\dd[3]{\vb y}\comm{\phi\qty(t,\vb x)}{\pi^\dagger\qty(t,\vb y)\pi\qty(t,\vb y)}\\
    \im\dot\phi\qty(t,\vb x)&=\int\dd[3]{\vb y}\pi^\dagger\qty(t,\vb y)\comm{\phi\qty(t,\vb x)}{\pi\qty(t,\vb y)}+\int\dd[3]{\vb y}\comm{\phi\qty(t,\vb x)}{\pi^\dagger\qty(t,\vb y)} \pi\qty(t,\vb y)\\
    \im\dot\phi\qty(t,\vb x)&=\int\dd[3]{\vb y}\pi^\dagger\qty(t,\vb y)\comm{\phi\qty(t,\vb x)}{\pi\qty(t,\vb y)}\\
    \im\dot\phi\qty(t,\vb x)&=\int\dd[3]{\vb y}\pi^\dagger\qty(t,\vb y)\im\delta^{\qty(3)}\qty(\vb x-\vb y)\\
    \im\dot\phi\qty(t,\vb x)&=\im\pi^\dagger\qty(t,\vb x)\Rightarrow \dot\phi\qty(t,\vb x)=\pi^\dagger\qty(t,\vb x)
\end{align}

Analogously is possible to show that,
\begin{align}
    \dot{\phi^\dagger} = \pi,\ \ \ \dot {\pi^\dagger} = \grad^2\phi-m^2\phi,\ \ \ \dot {\pi} = \grad^2\phi^\dagger-m^2\phi^\dagger
\end{align}
those are all the Heisenberg equations of motion for all the degrees of 
freedom. Notice that they can be put in the usual covariant form,
\begin{align}
    \dot {\pi^\dagger} &= \grad^2\phi-m^2\phi,\ \ \ \dot\phi\qty(t,\vb x)=\pi^\dagger\qty(t,\vb x)\\
    \ddot{\phi} &= \grad^2\phi-m^2\phi\Rightarrow -\Box\phi+m^2\phi = 0
\end{align}
which is a lot more familiar.

\probitem{}

% The canonical commutation relations are,
% \begin{align}
%     \comm{\phi\qty(t,\vb x)}{\pi\qty(t,\vb y)}&= \im\delta^{\qty(3)}\qty(\vb x-\vb y), \ \ \ \comm{\phi\qty(t,\vb x)}{\phi\qty(t,\vb y)}=\comm{\pi\qty(t,\vb x)}{\pi\qty(t,\vb y)}=0\\
%     \comm{\phi^\ast\qty(t,\vb x)}{\pi^\ast\qty(t,\vb y)}&= \im\delta^{\qty(3)}\qty(\vb x-\vb y), \ \ \ \comm{\phi^\ast\qty(t,\vb x)}{\phi^\ast\qty(t,\vb y)}=\comm{\pi^\ast\qty(t,\vb x)}{\pi^\ast\qty(t,\vb y)}=0\\
%     \comm{\phi^\ast\qty(t,\vb x)}{\phi\qty(t,\vb y)}&=\comm{\pi^\ast\qty(t,\vb x)}{\pi\qty(t,\vb y)}=0
% \end{align}

We have already stated the canonical commutation relations, and also computed 
the equations of motion. Now, we'll solve the equation of motion for $\phi$, 
this solution will have undetermined coefficients $a,b$, which also will have 
the interpretation of operators in the Hilbert space. We'll determine the algebra 
of these coefficients by ensuring $\phi$ satisfy the canonical commutation relations.

As we derived, our equation of motion is,
\begin{align}
    \qty(-\Box_x+m^2)\phi\qty(x)=0
\end{align}
which is linear. One of the easiest ways of solving linear PDE is 
to expand the field in it's inverse Fourier transform,
\begin{align}
    \phi\qty(x)&=\int\frac{\dd[4]{p}}{\qty(2\pi)^4}\tilde \phi\qty(p)\e^{\im x\cdot p}
\end{align}
substituting it back in the equation of motion,
\begin{align}
    0=\qty(-\Box_x+m^2)\phi\qty(x)&=\int\frac{\dd[4]{p}}{\qty(2\pi)^4}\tilde \phi\qty(p)\qty(-\Box_x+m^2)\e^{\im x\cdot p}\\
    0&=\int\frac{\dd[4]{p}}{\qty(2\pi)^4}\tilde \phi\qty(p)\qty(p^2+m^2)\e^{\im x\cdot p}
\end{align}
we obtained an integral condition of the Fourier modes $\tilde\phi$ which should be valid for every $x$. In fact, given an arbitrary function $\psi\qty(x)$, we can integrate over $x$,
\begin{align}
    0\int\dd[4]{x}\psi(x)&=\int\frac{\dd[4]{p}}{\qty(2\pi)^4}\tilde \phi\qty(p)\qty(p^2+m^2)\int\dd[4]{x}\e^{\im x\cdot p}\psi(x)\\
    0&=\int\frac{\dd[4]{p}}{\qty(2\pi)^4}\tilde \phi\qty(p)\qty(p^2+m^2)\tilde\psi(-p)
\end{align}
as $\psi$ is arbitrary\footnote{Actually $\psi$ has to be a test function, infinitely differentiable with continuous derivative and compact support. Nevertheless, the matter of whether $\phi$ do possesses an well defined Fourier transform is left untouched. 
There are developments into formalizing what kind of mathematical object $\phi$ --- an operator valued distribution ---, but we'll just assume 
everything works out.}, the only way to the result be zero is if $\tilde \phi\qty(p)\qty(p^2+m^2)$ is itself 
the zero distribution. Naively this would imply $\tilde \phi\qty(p) \equiv 0$, but, due to the distributional 
identity, \begin{align}x\delta(x)=0\end{align} the most general solution is $\tilde \phi\qty(p) \propto \delta\qty(p^2+m^2)$, as $\qty(p^2+m^2)\delta\qty(p^2+m^2)=0$. 
Our desired expression for $\tilde\phi$ is then given by $\delta\qty(p^2+m^2)$ multiplied by some coefficients, 
\begin{align}
    \tilde\phi\qty(p)&=A\qty(p)\delta\qty(p^2+m^2)
\end{align}
due to the Dirac delta, we're enforcing \begin{align}p^2 = -m^2\Rightarrow {p^0}^2 = \vb p^2 + m^2\Rightarrow p^0 =\pm\sqrt{ \vb p^2 + m^2}=\pm\omega_{\vb p}\end{align} 
thus, $A$ cannot have arbitrary arguments in the $p^0$ entry, otherwise the Dirac delta would enforce it 
to be zero. The only two non-zero contributions are $A\qty(\omega_{\vb p},\vb p)$ and $A\qty(-\omega_{\vb p},\vb p)$, 
\begin{align}
    \tilde\phi\qty(p)&=\qty(A\qty(\omega_{\vb p},\vb p)+A\qty(-\omega_{\vb p},\vb p))\delta\qty(p^2+m^2)
\end{align}

As the Dirac delta enforces $p^0=\pm\omega_{\vb p}$, it's not necessary to even label the $p^0$ in $A$, just 
to keep track of the sign. For this we call $A\qty(\omega_{\vb p},\vb p)=2\pi a\qty(\vb p)$ and 
$A\qty(-\omega_{\vb p},\vb p)=2\pi b^\dagger\qty(\vb p)$, and keep track of the $p^0$ sign by multiplying 
$a$ by $\theta\qty(p^0)$ and $b$ by $\theta\qty(-p^0)$,
\begin{align}
    \tilde\phi\qty(p)&=2\pi\delta\qty(p^2+m^2)\qty(\theta\qty(p^0)a\qty(\vb p)+\theta\qty(-p^0)b^\dagger\qty(\vb p))
\end{align}
of course the choice of normalization $2\pi$ is conventional, also the names of the coefficients $a, b^\dagger$. But, 
there are motifs behind this, which is a desired form of the final expression. Now, with the full expression of $\tilde\phi$, 
we do the Fourier transform to get $\phi$,
\begin{align}
    \phi\qty(x)&=\int\frac{\dd[4]{p}}{\qty(2\pi)^4}\tilde \phi\qty(p)\e^{\im x\cdot p}\\
    \phi\qty(x)&=\int\frac{\dd[4]{p}}{\qty(2\pi)^4}2\pi\delta\qty(p^2+m^2)\qty(\theta\qty(p^0)a\qty(\vb p)+\theta\qty(-p^0)b^\dagger\qty(\vb p))\e^{\im x\cdot p}
\end{align}
to simplify this we need to use some additional properties of the Dirac delta, namely,
\begin{align}
    \delta\qty(g(x)) = \sum\limits_{a}\frac{1}{\norm{g'\qty(x_a)}}\delta\qty(x-x_a),\ \ \ g(x_a) = 0
\end{align}
for our case $g(p^0) = -{p^0}^2+\omega_{\vb p}^2$, what implies,
\begin{align}
    \phi\qty(x)&=\int\frac{\dd[4]{p}}{\qty(2\pi)^32\omega_{\vb p}}\qty(\delta\qty(p^0-\omega_{\vb p})+\delta\qty(p^0+\omega_{\vb p}))\qty(\theta\qty(p^0)a\qty(\vb p)+\theta\qty(-p^0)b^\dagger\qty(\vb p))\e^{\im x\cdot p}\\
    \phi\qty(x)&=\int\frac{\dd[4]{p}}{\qty(2\pi)^32\omega_{\vb p}}\qty(\delta\qty(p^0-\omega_{\vb p})a\qty(\vb p)+\delta\qty(p^0+\omega_{\vb p})b^\dagger\qty(\vb p))\e^{\im x\cdot p}\\
    \phi\qty(x)&=\int\frac{\dd[3]{\vb p}}{\qty(2\pi)^32\omega_{\vb p}}a\qty(\vb p)\eval_{p^0=\omega_{\vb p}}+\int\frac{\dd[3]{\vb p}}{\qty(2\pi)^32\omega_{\vb p}}b^\dagger\qty(\vb p)\e^{\im x\cdot p}\eval_{p^0=-\omega_{\vb p}}
\end{align}
if we make the change of integration variables $\vb p\rightarrow-\vb p$ in the second integral, we can change the overall sign on the exponential to get,
\begin{align}
    \phi\qty(x)&=\int\frac{\dd[3]{\vb p}}{\qty(2\pi)^32\omega_{\vb p}}a\qty(\vb p)\eval_{p^0=\omega_{\vb p}}+\int\frac{\dd[3]{\vb p}}{\qty(2\pi)^32\omega_{\vb p}}b^\dagger\qty(\vb p)\e^{-\im x\cdot p}\eval_{p^0=\omega_{\vb p}}\\
    \phi\qty(x)&=\int\frac{\dd[3]{\vb p}}{\qty(2\pi)^32\omega_{\vb p}}\qty(a\qty(\vb p)\e^{\im x\cdot p}+b^\dagger\qty(\vb p)\e^{-\im x\cdot p})\eval_{p^0=\omega_{\vb p}}
\end{align}
if $\phi$ was a real field, we would have to impose $\phi^\dagger = \phi$, what would set $a = b$. But here we are working with a complex field, then $a$ and $b$ are independent.

Now we compute $\pi$ and make use of the commutation relations of the fields to obtain the commutation relations of the creation add 
annihilation operators. First,
\begin{align}
    \pi &= \dot{\phi^\dagger}\\
    \pi\qty(x)&=\int\frac{\dd[3]{\vb p}}{\qty(2\pi)^32\omega_{\vb p}}\im\omega_{\vb p}\qty(a^\dagger\qty(\vb p)\e^{-\im x\cdot p}-b\qty(\vb p)\e^{\im x\cdot p})\eval_{p^0=\omega_{\vb p}}
\end{align}
now, we do an inverse Fourier transform to isolate $a$,
\begin{align}
    \int\dd[3]{\vb x}e^{-\im x\cdot k}\phi\qty(x)&=\int\dd[3]{\vb x}e^{-\im x\cdot k}\int\frac{\dd[3]{\vb p}}{\qty(2\pi)^32\omega_{\vb p}}\qty(a\qty(\vb p)\e^{\im x\cdot p}+b^\dagger\qty(\vb p)\e^{-\im x\cdot p})\eval_{p^0=\omega_{\vb p}}\\
    \int\dd[3]{\vb x}e^{-\im x\cdot k}\phi\qty(x)&=\int\frac{\dd[3]{\vb p}}{\qty(2\pi)^32\omega_{\vb p}}\qty(\qty(2\pi)^3\delta^{\qty(3)}\qty(\vb k-\vb p)a\qty(\vb p)\e^{-\im x^0\qty( p^0-k^0)}+\qty(2\pi)^3\delta^{\qty(3)}\qty(\vb k+\vb p)b^\dagger\qty(\vb p)\e^{\im x^0\qty(p^0+k^0)})\eval_{p^0=\omega_{\vb p}}\\
    \int\dd[3]{\vb x}e^{-\im x\cdot k}\phi\qty(x)\eval_{k^0=\omega_{\vb k}}&=\frac{1}{2\omega_{\vb k}}\qty(a\qty(\vb k)+b^\dagger\qty(-\vb k)\e^{\im 2x^0k^0})\eval_{k^0=\omega_{\vb k}}
\end{align}
again but with $\pi^\dagger$,
\begin{align}
    \int\dd[3]{\vb x}e^{-\im x\cdot k}\pi^\dagger\qty(x)&=\int\dd[3]{\vb x}e^{-\im x\cdot k}\int\frac{\dd[3]{\vb p}}{\qty(2\pi)^32\omega_{\vb p}}\im\omega_{\vb p}\qty(-a\qty(\vb p)\e^{\im x\cdot p}+b^\dagger\qty(\vb p)\e^{-\im x\cdot p})\eval_{p^0=\omega_{\vb p}}\\
    \int\dd[3]{\vb x}e^{-\im x\cdot k}\pi^\dagger\qty(x)&=\im\int\frac{\dd[3]{\vb p}}{\qty(2\pi)^32}\qty(-\qty(2\pi)^3\delta^{\qty(3)}\qty(\vb k-\vb p)a\qty(\vb p)\e^{-\im x^0\qty( p^0-k^0)}+\qty(2\pi)^3\delta^{\qty(3)}\qty(\vb k+\vb p)b^\dagger\qty(\vb p)\e^{\im x^0\qty(p^0+k^0)})\eval_{p^0=\omega_{\vb p}}\\
    \int\dd[3]{\vb x}e^{-\im x\cdot k}\pi^\dagger\qty(x)\eval_{k^0=\omega_{\vb k}}&=\im\frac{1}{2}\qty(-a\qty(\vb k)+b^\dagger\qty(-\vb k)\e^{\im 2x^0k^0})\eval_{k^0=\omega_{\vb k}}
\end{align}

From these results, it's clear that,
\begin{align}
    \omega_{\vb k}\int\dd[3]{\vb x}e^{-\im x\cdot k}\eval_{k^0=\omega_{\vb k}}\phi\qty(t,\vb x)+\im\int\dd[3]{\vb x}e^{-\im x\cdot k}\eval_{k^0=\omega_{\vb k}}\pi^\dagger\qty(t,\vb x) = a\qty(\vb k)
\end{align}
with also the other possible combinations giving,
\begin{align}
    a^\dagger\qty(\vb k)&=\int\dd[3]{\vb x}e^{\im x\cdot k}\qty(\omega_{\vb k}\phi^\dagger\qty(t,\vb x)-\im\pi\qty(t,\vb x))\eval_{k^0=\omega_{\vb k}}\\
    b\qty(\vb k)&=\int\dd[3]{\vb x}e^{-\im x\cdot k}\qty(\omega_{\vb k}\phi^\dagger\qty(t,\vb x)+\im\pi\qty(t,\vb x))\eval_{k^0=\omega_{\vb k}}\\
    b^\dagger\qty(\vb k)&=\int\dd[3]{\vb x}e^{\im x\cdot k}\qty(\omega_{\vb k}\phi\qty(t,\vb x)-\im\pi^\dagger\qty(t,\vb x))\eval_{k^0=\omega_{\vb k}}
\end{align}

It is then straightforward to compute the commutation relations. As an example we compute one of the non-trivial,
\begin{align}
    \comm{a(\vb k)}{a^\dagger\qty(\vb p)} &=\int\dd[3]{\vb x}e^{-\im x\cdot k}\eval_{k^0=\omega_{\vb k}}\int\dd[3]{\vb y}e^{\im y\cdot p}\eval_{p^0=\omega_{\vb p}}\comm{\omega_{\vb k}\phi\qty(t,\vb x)+\im\pi^\dagger\qty(t,\vb x)}{\omega_{\vb p}\phi^\dagger\qty(t,\vb y)-\im\pi\qty(t,\vb y)}\\
    \comm{a(\vb k)}{a^\dagger\qty(\vb p)} &=\int\dd[3]{\vb x}e^{-\im x\cdot k}\eval_{k^0=\omega_{\vb k}}\int\dd[3]{\vb y}e^{\im y\cdot p}\eval_{p^0=\omega_{\vb p}}\im\qty(-\omega_{\vb k}\comm{\phi\qty(t,\vb x)}{\pi\qty(t,\vb y)}+\omega_{\vb p}\comm{\pi^\dagger\qty(t,\vb x)}{\phi^\dagger\qty(t,\vb y)})\\
    \comm{a(\vb k)}{a^\dagger\qty(\vb p)} &=\int\dd[3]{\vb x}e^{-\im x\cdot k}\eval_{k^0=\omega_{\vb k}}\int\dd[3]{\vb y}e^{\im y\cdot p}\eval_{p^0=\omega_{\vb p}}\im\qty(-\omega_{\vb k}\im\delta^{\qty(3)}\qty(\vb x-\vb y )-\omega_{\vb p}\im\delta^{\qty(3)}\qty(\vb x- \vb y))\\
    \comm{a(\vb k)}{a^\dagger\qty(\vb p)} &=\int\dd[3]{\vb x}e^{-\im x\cdot k}\eval_{k^0=\omega_{\vb k}}e^{\im x\cdot p}\eval_{p^0=\omega_{\vb p}}\qty(\omega_{\vb k}+\omega_{\vb p})\\
    \comm{a(\vb k)}{a^\dagger\qty(\vb p)} &=\qty(2\pi)^3\delta^{\qty(3)}\qty(\vb p -\vb k)\qty(\omega_{\vb k}+\omega_{\vb p})\\
    \comm{a(\vb k)}{a^\dagger\qty(\vb p)} &=\qty(2\pi)^3\omega_{\vb p }\delta^{\qty(3)}\qty(\vb p -\vb k)
\end{align}
analogously we can show,
\begin{align}
    \comm{a(\vb k)}{a\qty(\vb p)}&=\comm{a(\vb k)}{b^\dagger\qty(\vb p)}=\comm{a(\vb k)}{b\qty(\vb p)}=\comm{b(\vb k)}{b\qty(\vb p)}=0\\
    \comm{b(\vb k)}{b^\dagger\qty(\vb p)}&=\qty(2\pi)^3\omega_{\vb p }\delta^{\qty(3)}\qty(\vb p -\vb k)
\end{align}

These values of commutators for the modes of $\phi$ makes very clear that we have two decoupled particles, which are created by $a^\dagger,b^\dagger$ and annihilated by 
$a,b$. Both of them have positive energy and the same mass, but, as we'll see, they have opposite sign for the conserved charge of phase symmetry.

\probitem{}
\label{section:e}

These next computations are quite cumbersome and lengthy, despite this, the final results are very intuitive and clean. Let's start with the Hamiltonian which we already derived it's form in the preceding 
items,
\begin{align}
    H&=\int\dd[3]{\vb x}\qty[\pi^\ast\pi+\grad\phi^\ast\cdot\grad\phi +m^2\phi^\ast\phi]
\end{align}
next we substitute the mode expansion of $\phi$ and $\pi$,
\begin{align}
    H&=\int\dd[3]{\vb x}\int\frac{\dd[3]{\vb p}\dd[3]{\vb k}}{\qty(2\pi)^64\omega_{\vb p}\omega_{\vb k}}\left[\im\omega_{\vb k}\qty(-a\qty(\vb k)\e^{\im x\cdot k}+b^\dagger\qty(\vb k)\e^{-\im x\cdot k})\im\omega_{\vb p}\qty(-b\qty(\vb p)\e^{\im x\cdot p}+a^\dagger\qty(\vb p)\e^{-\im x\cdot p})\right.\nonumber\\
    &\quad\quad\quad+\grad\qty(b\qty(\vb k)\e^{\im x\cdot k}+a^\dagger\qty(\vb k)\e^{-\im x\cdot k})\cdot\grad\qty(a\qty(\vb p)\e^{\im x\cdot p}+b^\dagger\qty(\vb p)\e^{-\im x\cdot p})\nonumber\\
    &\quad\quad\quad\left.+m^2\qty(b\qty(\vb k)\e^{\im x\cdot k}+a^\dagger\qty(\vb k)\e^{-\im x\cdot k})\qty(a\qty(\vb p)\e^{\im x\cdot p}+b^\dagger\qty(\vb p)\e^{-\im x\cdot p})\right]\eval_{p^0=\omega_{\vb p}}\eval_{k^0=\omega_{\vb k}}\\
    H&=\int\dd[3]{\vb x}\int\frac{\dd[3]{\vb p}\dd[3]{\vb k}}{\qty(2\pi)^64\omega_{\vb p}\omega_{\vb k}}\left[\im\omega_{\vb k}\qty(-a\qty(\vb k)\e^{\im x\cdot k}+b^\dagger\qty(\vb k)\e^{-\im x\cdot k})\im\omega_{\vb p}\qty(-b\qty(\vb p)\e^{\im x\cdot p}+a^\dagger\qty(\vb p)\e^{-\im x\cdot p})\right.\nonumber\\
    &\quad\quad\quad+\im^2\vb p\cdot \vb k\qty(b\qty(\vb k)\e^{\im x\cdot k}-a^\dagger\qty(\vb k)\e^{-\im x\cdot k})\qty(a\qty(\vb p)\e^{\im x\cdot p}-b^\dagger\qty(\vb p)\e^{-\im x\cdot p})\nonumber\\
    &\quad\quad\quad\left.+m^2\qty(b\qty(\vb k)\e^{\im x\cdot k}+a^\dagger\qty(\vb k)\e^{-\im x\cdot k})\qty(a\qty(\vb p)\e^{\im x\cdot p}+b^\dagger\qty(\vb p)\e^{-\im x\cdot p})\right]\eval_{p^0=\omega_{\vb p}}\eval_{k^0=\omega_{\vb k}}
\end{align}
to ease our computation, we'll collect similar exponentials,
\begin{align}
    H&=\int\dd[3]{\vb x}\int\frac{\dd[3]{\vb p}\dd[3]{\vb k}}{\qty(2\pi)^64\omega_{\vb p}\omega_{\vb k}}\left[\e^{\im x\cdot\qty(k+p)}\qty(-\omega_{\vb k}\omega_{\vb p}a\qty(\vb k)b\qty(\vb p)-\vb k\cdot \vb p b\qty(\vb k)a\qty(\vb p)+m^2b\qty(\vb k)a\qty(\vb p))\right.\nonumber\\
    &\quad\quad\quad+\e^{\im x\cdot\qty(k-p)}\qty(\omega_{\vb k}\omega_{\vb p}a\qty(\vb k)a^\dagger\qty(\vb p)+\vb k\cdot \vb p b\qty(\vb k)b^\dagger\qty(\vb p)+m^2b\qty(\vb k)b^\dagger\qty(\vb p))\nonumber\\
    &\quad\quad\quad+\e^{\im x\cdot\qty(-k+p)}\qty(\omega_{\vb k}\omega_{\vb p}b^\dagger\qty(\vb k)b\qty(\vb p)+\vb k\cdot \vb p a^\dagger\qty(\vb k)a\qty(\vb p)+m^2a^\dagger\qty(\vb k)a\qty(\vb p))\nonumber\\
    &\quad\quad\quad\left.+\e^{-\im x\cdot\qty(k+p)}\qty(-\omega_{\vb k}\omega_{\vb p}b^\dagger\qty(\vb k)a^\dagger\qty(\vb p)-\vb k\cdot \vb p a^\dagger\qty(\vb k)b^\dagger\qty(\vb p)+m^2a^\dagger\qty(\vb k)b^\dagger\qty(\vb p))\right]\eval_{p^0=\omega_{\vb p}}\eval_{k^0=\omega_{\vb k}}
\end{align}
now we perform the $\dd[3]{\vb x}$ integral which will give Dirac deltas, and also perform the $\dd[3]{\vb k} $ integrals over the deltas,
\begin{align}
    H&=\int\frac{\dd[3]{\vb p}}{\qty(2\pi)^34\omega_{\vb p}^2}\left[\e^{-\im x^02\omega_{\vb p}}\qty(-\omega_{\vb p}\omega_{\vb p}a\qty(-\vb p)b\qty(\vb p)+\vb p^2 b\qty(-\vb p)a\qty(\vb p)+m^2b\qty(-\vb p)a\qty(\vb p))\right.\nonumber\\
    &\quad\quad\quad+\qty(\omega_{\vb p}\omega_{\vb p}a\qty(\vb p)a^\dagger\qty(\vb p)+ \vb p^2 b\qty(\vb p)b^\dagger\qty(\vb p)+m^2b\qty(\vb p)b^\dagger\qty(\vb p))\nonumber\\
    &\quad\quad\quad+\qty(\omega_{\vb p}\omega_{\vb p}b^\dagger\qty(\vb p)b\qty(\vb p)+\vb p^2 a^\dagger\qty(\vb p)a\qty(\vb p)+m^2a^\dagger\qty(\vb p)a\qty(\vb p))\nonumber\\
    &\quad\quad\quad\left.+\e^{\im x^02\omega_{\vb p}}\qty(-\omega_{\vb p}\omega_{\vb p}b^\dagger\qty(-\vb p)a^\dagger\qty(\vb p)+\vb p^2 a^\dagger\qty(-\vb p)b^\dagger\qty(\vb p)+m^2a^\dagger\qty(-\vb p)b^\dagger\qty(\vb p))\right]
\end{align}
the first and last line are identically zero, because $a,b$ commute, as well as $a^\dagger,b^\dagger$, and $\omega_{\vb p}^2=\vb p^2+m^2$. We also use this last equality to collect the $\vb p^2,m^2$ terms in the second and third lines,
\begin{align}
    H&=\int\frac{\dd[3]{\vb p}}{\qty(2\pi)^34\omega_{\vb p}^2}\left[\omega_{\vb p}^2a\qty(\vb p)a^\dagger\qty(\vb p)+ \omega_{\vb p}^2 b\qty(\vb p)b^\dagger\qty(\vb p)+\omega_{\vb p}^2b^\dagger\qty(\vb p)b\qty(\vb p)+\omega_{\vb p}^2 a^\dagger\qty(\vb p)a\qty(\vb p)\right]
\end{align}
making use of the commutation relations we exchange the order of the operators to all have the annihilation ones to the right,
\begin{align}
    H&=\int\frac{\dd[3]{\vb p}}{\qty(2\pi)^34\omega_{\vb p}}\omega_{\vb p}\left[2\qty(2\pi)^32\omega_{\vb p}\delta^{\qty(3)}\qty(\vb 0)+2b^\dagger\qty(\vb p)b\qty(\vb p)+ 2a^\dagger\qty(\vb p)a\qty(\vb p)\right]\\
    H&=\int\frac{\dd[3]{\vb p}}{\qty(2\pi)^32\omega_{\vb p}}\omega_{\vb p}\left[b^\dagger\qty(\vb p)b\qty(\vb p)+ a^\dagger\qty(\vb p)a\qty(\vb p)\right]+2\delta^{\qty(3)}\qty(\vb 0)\int\dd[3]{\vb p}\frac12\omega_{\vb p}
\end{align}

Let's comment a little about this last expression of the Hamiltonian operator. The combination $aa^\dagger$, or $bb^\dagger$, is analogous from usual quantum mechanics, this is the occupation number operator for $a$, or $b$. 
That is, it counts how many of $a(b)$-type particles are present in a state. We're multiplying it by $\omega_{\vb p}$, and integrating over $\frac{\dd[3]{\vb p}}{\qty(2\pi)^32\omega_{\vb p}}$. In other words, 
we count the number of particles with a certain energy $\omega_{\vb p} $ by applying $a^\dagger(\vb p)a\qty(\vb p)$, we multiply the number of particles with energy $\omega_{\vb p}$ by the energy itself, and lastly, 
we sum over all possible energies with the measure\footnote{This is conventional, it depends on normalization choices. In this choice, the integration measure is Lorentz invariant and $a,b,a^\dagger,b^\dagger$ transform 
as scalars.} $\frac{\dd[3]{\vb p}}{\qty(2\pi)^32\omega_{\vb p}}$. The final answer is the total energy!

Well, kind of. Notice that we have a constant term $2\delta^{\qty(3)}\qty(\vb 0)\int\dd[3]{\vb k}\omega_{\vb k}$, which is infinite. Being infinite is not really a problem as 
we are considering the total energy of an infinite spatial volume, as $\delta^{(3)}\qty(\vb 0)=\frac{V}{\qty(2\pi)^3}$, where $V$ is 
the spatial volume. We could consider the energy density,
\begin{align}
    \frac1V2\delta^{\qty(3)}\qty(\vb 0)\int\dd[3]{\vb p}\frac12\omega_{\vb p}=\frac{2}{\qty(2\pi)^3}\int\dd[3]{\vb p}\frac12\omega_{\vb p}
\end{align}
which is still divergent, this is a ``problem''. There are many ways to deal with this. A common argument is that we can only measure differences in energy between states, and as this divergent term is a constant, 
it would drop out in differences. Despite being true, it isn't that satisfactory answer, happily there are more interesting answers. One of the main point is to look at the origin of this constant term, it appears 
as we tried to change the ordering of the operators in the expression of the Hamiltonian, but notice, we inherited the expression for the quantum Hamiltonian from the classical expression, by reinterpreting the 
classical commuting phase space variables as quantum non-commuting operators. In the classical expression we could have chosen any ordering of our liking, as the dynamical variables commute, this is alarming! 
If we could have chosen any ordering in the classical expression, from which our quantum expression comes, that means we could have started the computation of the quantum expression of the Hamiltonian from 
any ordering! There is an ordering ambiguity when quantizing a theory! At least this is the way the problem presents itself in the canonical quantization. There are two other main ways of dealing with this. 
One is to define a ordering prescription for all operators, to remove the ambiguity. We call it the normal ordering. Given an operator $\mathcal O$ which is a function of creation and annihilation 
operators, we define it's normal ordering as $\cnord{\mathcal O}$. Operationally we move, by hand, all annihilation operators to the right of all creation operators, with the additional property $\cnord{\mathbbm1}=0$. 
This is equivalent to defining $\cnord{\mathcal O} \coloneq \mathcal O-\mel{0}{\mathcal O}{0}$. If we apply to the Hamiltonian,
\begin{align}
    \cnord{H}&=\int\frac{\dd[3]{\vb p}}{\qty(2\pi)^32\omega_{\vb p}}\omega_{\vb p}\left[b^\dagger\qty(\vb p)b\qty(\vb p)+ a^\dagger\qty(\vb p)a\qty(\vb p)\right]
\end{align}

The other way of dealing with this is to notice from the definition of the Hamiltonian,
\begin{align}
    \mathcal H  &= \pi^I\dot{ \phi_I} - \mathcal L
\end{align}
that we can include an arbitrary constant term in $\mathcal L$ which will not affect in any way the equations of motion, but will shift the Hamiltonian by a constant term. Say that we include a constant $\Omega$ in the 
Lagrangian, $\mathcal L \rightarrow \mathcal L +\Omega$, then, $\mathcal H\rightarrow \mathcal H-\Omega$, and,
\begin{align}
    H&=\int\frac{\dd[3]{\vb p}}{\qty(2\pi)^32\omega_{\vb p}}\omega_{\vb p}\left[b^\dagger\qty(\vb p)b\qty(\vb p)+ a^\dagger\qty(\vb p)a\qty(\vb p)\right]+2\delta^{\qty(3)}\qty(\vb 0)\int\dd[3]{\vb p}\frac12\omega_{\vb p}-\int\dd[3]{\vb x}\Omega\\
    H&=\int\frac{\dd[3]{\vb p}}{\qty(2\pi)^32\omega_{\vb p}}\omega_{\vb p}\left[b^\dagger\qty(\vb p)b\qty(\vb p)+ a^\dagger\qty(\vb p)a\qty(\vb p)\right]+2\int\dd[3]{\vb x}\int\frac{\dd[3]{\vb p}}{\qty(2\pi)^3}\frac12\omega_{\vb p}-\int\dd[3]{\vb x}\Omega
\end{align}
by choosing, 
\begin{align}
    \Omega&=2\int\frac{\dd[3]{\vb p}}{\qty(2\pi)^3}\frac12\omega_{\vb p}
\end{align}
we arrive at a finite energy. 

This procedure is in it's essence a renormalization of the cosmological constant, which is in practice equivalent to choose a operator ordering.

Now we follow with the computation of the linear momenta, the way to go is from the Noether charge associated with translations, as a reminder,
\begin{align}
    \tensor{T}{^\mu^\nu}&=-\pdv{\mathcal L}{\partial_\mu\phi_I}\partial^\nu\phi_I+g^{\mu\nu}\mathcal L\\
    P^\nu&= \int\dd[3]{\vb x}T^{0\nu}
\end{align}
of course this gives the Hamiltonian as $P^0=H$. But let's focus on the linear momentum, $\nu=i$,
\begin{align}
    P^i&= -\int\dd[3]{\vb x}\pi^I\partial^i\phi_I =  -\int\dd[3]{\vb x}\qty(\pi\partial^i\phi+\pi^\dagger\partial^i\phi^\dagger)
\end{align}
now substitute the mode expansion,
\begin{align}
    P^i&=-\int\dd[3]{\vb x}\int\frac{\dd[3]{\vb k}\dd[3]{\vb p}}{\qty(2\pi)^64\omega_{\vb k}\omega_{\vb p}}\im\omega_{\vb k}\qty(a^\dagger\qty(\vb k)\e^{-\im x\cdot k}-b\qty(\vb k)\e^{\im x\cdot k})\eval_{k^0=\omega_{\vb k}}\partial^i\qty(a\qty(\vb p)\e^{\im x\cdot p}+b^\dagger\qty(\vb p)\e^{-\im x\cdot p})\eval_{p^0=\omega_{\vb p}}\nonumber\\
    &\quad\quad\quad-\int\dd[3]{\vb x}\int\frac{\dd[3]{\vb k}\dd[3]{\vb p}}{\qty(2\pi)^64\omega_{\vb k}\omega_{\vb p}}\qty(-\im\omega_{\vb k})\qty(a\qty(\vb k)\e^{\im x\cdot k}-b^\dagger\qty(\vb k)\e^{-\im x\cdot k})\eval_{k^0=\omega_{\vb k}}\partial^i\qty(a^\dagger\qty(\vb p)\e^{-\im x\cdot p}+b\qty(\vb p)\e^{\im x\cdot p})\eval_{p^0=\omega_{\vb p}}\\
    P^i&=-\int\dd[3]{\vb x}\int\frac{\dd[3]{\vb k}\dd[3]{\vb p}}{\qty(2\pi)^64\omega_{\vb p}}\im\qty(a^\dagger\qty(\vb k)\e^{-\im x\cdot k}-b\qty(\vb k)\e^{\im x\cdot k})\eval_{k^0=\omega_{\vb k}}\im p^i\qty(a\qty(\vb p)\e^{\im x\cdot p}-b^\dagger\qty(\vb p)\e^{-\im x\cdot p})\eval_{p^0=\omega_{\vb p}}\nonumber\\
    &\quad\quad\quad-\int\dd[3]{\vb x}\int\frac{\dd[3]{\vb k}\dd[3]{\vb p}}{\qty(2\pi)^64\omega_{\vb p}}\qty(-\im)\qty(a\qty(\vb k)\e^{\im x\cdot k}-b^\dagger\qty(\vb k)\e^{-\im x\cdot k})\eval_{k^0=\omega_{\vb k}}\im p^i\qty(-a^\dagger\qty(\vb p)\e^{-\im x\cdot p}+b\qty(\vb p)\e^{\im x\cdot p})\eval_{p^0=\omega_{\vb p}}\\
    P^i&=\int\dd[3]{\vb x}\int\frac{\dd[3]{\vb k}\dd[3]{\vb p}}{\qty(2\pi)^64\omega_{\vb p}}\qty(a^\dagger\qty(\vb k)\e^{-\im x\cdot k}-b\qty(\vb k)\e^{\im x\cdot k})\eval_{k^0=\omega_{\vb k}} p^i\qty(a\qty(\vb p)\e^{\im x\cdot p}-b^\dagger\qty(\vb p)\e^{-\im x\cdot p})\eval_{p^0=\omega_{\vb p}}\nonumber\\
    &\quad\quad\quad-\int\dd[3]{\vb x}\int\frac{\dd[3]{\vb k}\dd[3]{\vb p}}{\qty(2\pi)^64\omega_{\vb p}}\qty(a\qty(\vb k)\e^{\im x\cdot k}-b^\dagger\qty(\vb k)\e^{-\im x\cdot k})\eval_{k^0=\omega_{\vb k}} p^i\qty(-a^\dagger\qty(\vb p)\e^{-\im x\cdot p}+b\qty(\vb p)\e^{\im x\cdot p})\eval_{p^0=\omega_{\vb p}}
\end{align}
collecting the similar exponentials,
\begin{align}
    P^i&=\int\dd[3]{\vb x}\int\frac{\dd[3]{\vb k}\dd[3]{\vb p}}{\qty(2\pi)^64\omega_{\vb p}}p^i\left[\e^{\im x\cdot\qty(p-k)}\qty(a^\dagger\qty(\vb k)a\qty(\vb p)+b^\dagger\qty(\vb k)b\qty(\vb p))\right.\nonumber\\
    &\quad\quad\quad-\e^{-\im x\cdot\qty(p+k)}\qty(a^\dagger\qty(\vb k)b^\dagger\qty(\vb p)+b^\dagger\qty(\vb k)a^\dagger\qty(\vb p))\nonumber\\
    &\quad\quad\quad-\e^{\im x\cdot\qty(p+k)}\qty(b\qty(\vb k)a\qty(\vb p)+a\qty(\vb k)b\qty(\vb p))\nonumber\\
    &\quad\quad\quad\left.+\e^{\im x\cdot\qty(k-p)}\qty(b\qty(\vb k)b^\dagger\qty(\vb p)+a\qty(\vb k)a^\dagger\qty(\vb p))\right]\eval_{k^0=\omega_{\vb k}}\eval_{p^0=\omega_{\vb p}}
\end{align}
doing the $\dd[3]{\vb x}$ integral and integrating the Dirac deltas in $\dd[3]{\vb k}$,
\begin{align}
    P^i&=\int\frac{\dd[3]{\vb p}}{\qty(2\pi)^64\omega_{\vb p}}p^i\left[\qty(2\pi)^3\qty(a^\dagger\qty(\vb p)a\qty(\vb p)+b^\dagger\qty(\vb p)b\qty(\vb p))\right.\nonumber\\
    &\quad\quad\quad-\e^{2\im x^0p^0}\qty(2\pi)^3\qty(a^\dagger\qty(-\vb p)b^\dagger\qty(\vb p)+b^\dagger\qty(-\vb p)a^\dagger\qty(\vb p))\nonumber\\
    &\quad\quad\quad-\e^{-2\im x^0p^0}\qty(2\pi)^3\qty(b\qty(-\vb p)a\qty(\vb p)+a\qty(-\vb p)b\qty(\vb p))\nonumber\\
    &\quad\quad\quad\left.+\qty(2\pi)^3\qty(b\qty(\vb p)b^\dagger\qty(\vb p)+a\qty(\vb p)a^\dagger\qty(\vb p))\right]\eval_{p^0=\omega_{\vb k}}
\end{align}
the second and third lines are zero, this is so because $a,b$ commute, as well as $a^\dagger,b^\dagger$, then, we can exchange the integration variables in the second term of each of these lines as $\vb p\rightarrow -\vb p$. 
Due to the overall $p^i$ factor, there is a global minus sign in this transformation, then,
\begin{align}
    P^i&=\int\frac{\dd[3]{\vb p}}{\qty(2\pi)^34\omega_{\vb p}}p^i\left[\qty(a^\dagger\qty(\vb p)a\qty(\vb p)+b^\dagger\qty(\vb p)b\qty(\vb p))\right.\nonumber\\
    &\quad\quad\quad-\e^{2\im x^0p^0}\qty(a^\dagger\qty(-\vb p)b^\dagger\qty(\vb p)-a^\dagger\qty(-\vb p)b^\dagger\qty(\vb p))\nonumber\\
    &\quad\quad\quad-\e^{-2\im x^0p^0}\qty(b\qty(-\vb p)a\qty(\vb p)-b\qty(-\vb p)a\qty(\vb p))\nonumber\\
    &\quad\quad\quad\left.+\qty(b\qty(\vb p)b^\dagger\qty(\vb p)+a\qty(\vb p)a^\dagger\qty(\vb p))\right]\eval_{p^0=\omega_{\vb k}}\\
    P^i&=\int\frac{\dd[3]{\vb p}}{\qty(2\pi)^34\omega_{\vb p}}p^i\qty(a^\dagger\qty(\vb p)a\qty(\vb p)+b^\dagger\qty(\vb p)b\qty(\vb p)+b\qty(\vb p)b^\dagger\qty(\vb p)+a\qty(\vb p)a^\dagger\qty(\vb p))
\end{align}
lastly, using the commutation relations to change the operators ordering,
\begin{align}
    P^i&=\int\frac{\dd[3]{\vb p}}{\qty(2\pi)^34\omega_{\vb p}}p^i\qty(2a^\dagger\qty(\vb p)a\qty(\vb p)+2b^\dagger\qty(\vb p)b\qty(\vb p)+2\qty(2\pi)^3\delta^{\qty(3)}\qty(\vb 0)2\omega_{\vb p})\\
    P^i&=\int\frac{\dd[3]{\vb p}}{\qty(2\pi)^32\omega_{\vb p}}p^i\qty(a^\dagger\qty(\vb p)a\qty(\vb p)+b^\dagger\qty(\vb p)b\qty(\vb p))+2\delta^{\qty(3)}\qty(\vb 0)\int\dd[3]{\vb p}\frac12p^i
\end{align}

The interpretation of this final result is the same, the first term count the number of particles of $a(b)$-type with $\vb p$ momentum and multiply it by the value of the momentum, after, it sums over all possible momentum 
with the Lorentz invariant measure. There is also an additional constant infinite momentum term, we can give the same argument, we only measure momentum differences between particles, and also the normal ordering procedure 
can be applied.

Lastly we do the same substitution for the conserved charge,
\begin{align}
    Q&=-\im\int\dd[3]{\vb x}\qty(\pi\phi-\pi^\dagger\phi^\dagger)\\
    Q&=-\im\int\dd[3]{\vb x}\int\frac{\dd[3]{\vb k}\dd[3]{\vb p}}{\qty(2\pi)^64\omega_{\vb k}\omega_{\vb p}}\im\omega_{\vb k}\qty(a^\dagger\qty(\vb k)\e^{-\im x\cdot k}-b\qty(\vb k)\e^{\im x\cdot k})\eval_{k^0=\omega_{\vb k}}\qty(a\qty(\vb p)\e^{\im x\cdot p}+b^\dagger\qty(\vb p)\e^{-\im x\cdot p})\eval_{p^0=\omega_{\vb p}}\nonumber\\
    &\quad\quad\quad-\im\int\dd[3]{\vb x}\int\frac{\dd[3]{\vb k}\dd[3]{\vb p}}{\qty(2\pi)^64\omega_{\vb k}\omega_{\vb p}}\qty(-\im\omega_{\vb k})\qty(a\qty(\vb k)\e^{\im x\cdot k}-b^\dagger\qty(\vb k)\e^{-\im x\cdot k})\eval_{k^0=\omega_{\vb k}}\qty(a^\dagger\qty(\vb p)\e^{-\im x\cdot p}+b\qty(\vb p)\e^{\im x\cdot p})\eval_{p^0=\omega_{\vb p}}
\end{align}
collecting the similar exponentials,
\begin{align}
    Q&=\int\dd[3]{\vb x}\int\frac{\dd[3]{\vb k}\dd[3]{\vb p}}{\qty(2\pi)^64\omega_{\vb p}}\left[\e^{\im x\cdot\qty(p-k)}\qty(a^\dagger\qty(\vb k)a\qty(\vb p)-b^\dagger\qty(\vb k)b\qty(\vb p))\right.\nonumber\\
    &\quad\quad\quad+\e^{-\im x\cdot\qty(k+p)}\qty(a^\dagger\qty(\vb k)b^\dagger\qty(\vb p)-b^\dagger\qty(\vb k)a^\dagger\qty(\vb p))\nonumber\\
    &\quad\quad\quad+\e^{\im x\cdot\qty(k+p)}\qty(-b\qty(\vb k)a\qty(\vb p)+a\qty(\vb k)b\qty(\vb p))\nonumber\\
    &\quad\quad\quad\left.+\e^{\im x \cdot\qty(k-p)}\qty(-b\qty(\vb k)b^\dagger\qty(\vb p)+a\qty(\vb k)a^\dagger\qty(\vb p))\right]\eval_{k^0=\omega_{\vb k}}\eval_{p^0=\omega_{\vb p}}
\end{align}
doing the $\dd[3]{\vb x}$ integral and integrating the Dirac deltas with $\dd[3]{\vb k}$,
\begin{align}
    Q&=\int\frac{\dd[3]{\vb p}}{\qty(2\pi)^64\omega_{\vb p}}\left[\qty(2\pi)^3\qty(a^\dagger\qty(\vb p)a\qty(\vb p)-b^\dagger\qty(\vb p)b\qty(\vb p))\right.\nonumber\\
    &\quad\quad\quad+\qty(2\pi)^3\e^{\im 2x^0p^0}\qty(a^\dagger\qty(-\vb p)b^\dagger\qty(\vb p)-b^\dagger\qty(-\vb p)a^\dagger\qty(\vb p))\nonumber\\
    &\quad\quad\quad+\qty(2\pi)^3\e^{-\im2 x^0p^0}\qty(-b\qty(-\vb p)a\qty(\vb p)+a\qty(-\vb p)b\qty(\vb p))\nonumber\\
    &\quad\quad\quad\left.+\qty(2\pi)^3\qty(-b\qty(\vb k)b^\dagger\qty(\vb p)+a\qty(\vb k)a^\dagger\qty(\vb p))\right]\eval_{p^0=\omega_{\vb p}}
\end{align}
the second and third lines are zero, this can be seen by performing a change of integration variables in the second term of each, $\vb p\rightarrow -\vb p$, and also because $a,b$ and $a^\dagger,b^\dagger$ commute,
\begin{align}
    Q&=\int\frac{\dd[3]{\vb p}}{\qty(2\pi)^34\omega_{\vb p}}\qty(a^\dagger\qty(\vb p)a\qty(\vb p)-b^\dagger\qty(\vb p)b\qty(\vb p)-b\qty(\vb k)b^\dagger\qty(\vb p)+a\qty(\vb k)a^\dagger\qty(\vb p))
\end{align}
using the commutation relations to change the operator ordering,
\begin{align}
    Q&=\int\frac{\dd[3]{\vb p}}{\qty(2\pi)^34\omega_{\vb p}}\qty(2a^\dagger\qty(\vb p)a\qty(\vb p)-2b^\dagger\qty(\vb p)b\qty(\vb p)-\qty(2\pi)^32\omega_{\vb p}\delta^{\qty(3)}\qty(\vb 0)+\qty(2\pi)^32\omega_{\vb p}\delta^{\qty(3)}\qty(\vb 0))\\
    Q&=\int\frac{\dd[3]{\vb p}}{\qty(2\pi)^32\omega_{\vb p}}\qty(a^\dagger\qty(\vb p)a\qty(\vb p)-b^\dagger\qty(\vb p)b\qty(\vb p))
\end{align}

The interpretation of this conserved charge is easy o read, $Q$ counts the number of $a$ particles minus the number of $b$ particles. In other words, we could say the $a$ particles have charge $+1$ while the $b$ particles 
have $-1$ charge. Of course this is entirely conventional, as $Q$ is conserved we could have defined $-Q$ as the physical charge, which would assign a different sign for the charge of the particles. The meaningful 
physical take is that $a$ and $b$ have same modulus opposite sign charges, that is, the number of $a$ particles minus the number of $b$ particles is constant in time evolution, in order to have more 
of $a(b)$ particles is necessary to create more $b(a)$ particles. 

\probitem{}

This last exercise is to compute the Feynman propagator for the complex scalar field as the vacuum expectation value of a time-ordered 2 field product. The main form of the propagator is 
highly similar between various theories, this is due to unitarity, locality, causality, Poincaré invariance, positive energy and mass being very strong assumptions, we have very little freedom 
in constructing propagators without violating some of these hypothesis. Manly due to this the propagator of the complex scalar field is identical to the real scalar field.

The computation begins with the definition of the time ordered product, we assume $x-y$ is time-like, otherwise there might be some ambiguity on the value of the Heaviside theta,
\begin{align}
    \Delta_{\textnormal{F}}\qty(x,y) &= \mel{0}{\textnormal{T}\qty{\phi\qty(x)\phi^\dagger\qty(y)}}{0}\\
    \Delta_{\textnormal{F}}\qty(x,y) &= \theta\qty(x^0-y^0)\mel{0}{\phi\qty(x)\phi^\dagger\qty(y)}{0}+\theta\qty(y^0-x^0)\mel{0}{\phi^\dagger\qty(y)\phi\qty(x)}{0}
\end{align}
now we proceed using the modes expansion,
\begin{align}
    \Delta_{\textnormal{F}}\qty(x,y) &= \int\frac{\dd[3]{\vb p}\dd[3]{\vb k}}{\qty(2\pi)^64\omega_{\vb p}\omega_{\vb k}}\left[\theta\qty(x^0-y^0)\mel{0}{\qty(a\qty(\vb p)\e^{\im x\cdot p}+b^\dagger\qty(\vb p)\e^{-\im x\cdot p})\qty(a^\dagger\qty(\vb k)\e^{-\im y\cdot k}+b\qty(\vb k)\e^{\im y\cdot p})}{0}\right.\nonumber\\
    &\quad\quad\quad\left.+\theta\qty(y^0-x^0)\mel{0}{\qty(b\qty(\vb k)\e^{\im y\cdot k}+a^\dagger\qty(\vb k)\e^{-\im y\cdot k})\qty(b^\dagger\qty(\vb p)\e^{-\im x\cdot p}+a\qty(\vb p)\e^{\im x\cdot p})}{0}\right]\eval_{p^0=\omega_{\vb p}}\eval_{k^0=\omega_{\vb k}}
\end{align}
we know use the property that the annihilation operators annihilate the vacuum, $a\qty(\vb k)\ket 0=b\qty(\vb k)\ket 0 =0$, and that the creation operators acting to the left annihilate the bra vacuum, $\bra 0 a^\dagger\qty(\vb k)=\bra 0 b^\dagger\qty(\vb k)=0$,
\begin{align}
    \Delta_{\textnormal{F}}\qty(x,y) &= \int\frac{\dd[3]{\vb p}\dd[3]{\vb k}}{\qty(2\pi)^64\omega_{\vb p}\omega_{\vb k}}\left[\theta\qty(x^0-y^0)\mel{0}{a\qty(\vb p)\e^{\im x\cdot p}a^\dagger\qty(\vb k)\e^{-\im y\cdot k}}{0}\right.\nonumber\\
    &\quad\quad\quad\left.+\theta\qty(y^0-x^0)\mel{0}{b\qty(\vb k)\e^{\im y\cdot k}b^\dagger\qty(\vb p)\e^{-\im x\cdot p}}{0}\right]\eval_{p^0=\omega_{\vb p}}\eval_{k^0=\omega_{\vb k}}
\end{align}
now we use the commutation relations to change the operators ordering, this will produce terms as $a^\dagger a$ and $b^\dagger b$, which are zero because they act on the vacuum, then,
\begin{align}
    \Delta_{\textnormal{F}}\qty(x,y) &= \int\frac{\dd[3]{\vb p}\dd[3]{\vb k}}{\qty(2\pi)^64\omega_{\vb p}\omega_{\vb k}}\left[\theta\qty(x^0-y^0)\e^{-\im y\cdot k+\im x\cdot p}\qty(2\pi)^32\omega_{\vb p}\delta^{\qty(3)}\qty(\vb k-\vb p)\braket{0}{0}\right.\nonumber\\
    &\quad\quad\quad\left.+\theta\qty(y^0-x^0)\e^{-\im x\cdot p+\im y\cdot k}\qty(2\pi)^32\omega_{\vb p}\delta^{\qty(3)}\qty(\vb k-\vb p)\braket{0}{0}\right]\eval_{p^0=\omega_{\vb p}}\eval_{k^0=\omega_{\vb k}}
\end{align}
next we integrate the Dirac deltas with $\dd[3]{\vb k}$,
\begin{align}
    \Delta_{\textnormal{F}}\qty(x,y) &= \int\frac{\dd[3]{\vb p}}{\qty(2\pi)^32\omega_{\vb p}}\left[\theta\qty(x^0-y^0)\e^{\im \qty(x-y)\cdot p}+\theta\qty(y^0-x^0)\e^{-\im \qty(x-y)\cdot p}\right]\eval_{p^0=\omega_{\vb p}}\\
    \Delta_{\textnormal{F}}\qty(x,y) &= \int\frac{\dd[3]{\vb p}}{\qty(2\pi)^32\omega_{\vb p}}\e^{\im \qty(\vb x-\vb y)\cdot \vb p}\left(\theta\qty(x^0-y^0)\e^{-\im \qty(x^0-y^0)\omega_{\vb p}}+\theta\qty(y^0-x^0)\e^{\im \qty(x^0-y^0)\omega_{\vb p}}\right)
\end{align}
where in the last line we separated the exponentials in the spatial and temporal pieces, this will ease the next step.

For this we'll need an additional piece of information, the integral representation of the Heaviside function, which is derived in \cref{appendix:heaviside},
\begin{align}
    \theta\qty(x) = -\frac{1}{2\pi\im}\lim\limits_{\epsilon\rightarrow 0^+}\int\limits_{-\infty}^{+\infty}\dd{p}\frac{\e^{-\im xp}}{p+\im\epsilon}
\end{align}
in particular we'll use the following choice of variables, to keep the process clean,
\begin{align}
    \theta\qty(x^0-y^0) = -\frac{1}{2\pi\im}\lim\limits_{\epsilon\rightarrow 0^+}\int\limits_{-\infty}^{+\infty}\dd{p^0}\frac{\e^{-\im \qty(x^0-y^0)p^0}}{p^0+\im\epsilon}
\end{align}
which now we substitute in our expression of the Feynman propagator,
\begin{align}
    \Delta_{\textnormal{F}}\qty(x,y) &=-\frac{1}{2\pi\im}\lim\limits_{\epsilon\rightarrow 0^+} \int\frac{\dd[3]{\vb p}}{\qty(2\pi)^32\omega_{\vb p}}\e^{\im \qty(\vb x-\vb y)\cdot \vb p}\int\dd{p^0}\frac{1}{p^0+\im\epsilon}\left(\e^{-\im \qty(x^0-y^0)p^0}\e^{-\im \qty(x^0-y^0)\omega_{\vb p}}+\e^{\im \qty(x^0-y^0)p^0}\e^{\im \qty(x^0-y^0)\omega_{\vb p}}\right)\\
    \Delta_{\textnormal{F}}\qty(x,y) &=-\frac{1}{2\pi\im}\lim\limits_{\epsilon\rightarrow 0^+} \int\frac{\dd[3]{\vb p}}{\qty(2\pi)^32\omega_{\vb p}}\e^{\im \qty(\vb x-\vb y)\cdot \vb p}\int\dd{p^0}\frac{1}{p^0+\im\epsilon}\left(\e^{-\im \qty(x^0-y^0)\qty(p^0+\omega_{\vb p})}+\e^{\im \qty(x^0-y^0)\qty(p^0+\omega_{\vb p})}\right)
\end{align}
we make a change of variables in the $p^0$ integral as $p^0\rightarrow p^0-\omega_{\vb p}$,
\begin{align}
    \Delta_{\textnormal{F}}\qty(x,y) &=-\frac{1}{2\pi\im}\lim\limits_{\epsilon\rightarrow 0^+} \int\frac{\dd[3]{\vb p}}{\qty(2\pi)^32\omega_{\vb p}}\e^{\im \qty(\vb x-\vb y)\cdot \vb p}\int\dd{p^0}\frac{1}{p^0-\omega_{\vb p}+\im\epsilon}\left(\e^{-\im \qty(x^0-y^0)p^0}+\e^{\im \qty(x^0-y^0)p^0}\right)
\end{align}
and now change the sign of $p^0$ in the second term as $p^0\rightarrow -p^0$,
\begin{align}
    \Delta_{\textnormal{F}}\qty(x,y) &=-\frac{1}{2\pi\im}\lim\limits_{\epsilon\rightarrow 0^+} \int\frac{\dd[3]{\vb p}}{\qty(2\pi)^32\omega_{\vb p}}\e^{\im \qty(\vb x-\vb y)\cdot \vb p}\int\dd{p^0}\e^{-\im \qty(x^0-y^0)p^0}\left(\frac{1}{p^0-\omega_{\vb p}+\im\epsilon}+\frac{1}{-p^0-\omega_{\vb p}+\im\epsilon}\right)
\end{align}
remains to collect both fractions, and to merge the two exponentials as a single four product,
\begin{align}
    \Delta_{\textnormal{F}}\qty(x,y) &=-\frac{1}{\im}\lim\limits_{\epsilon\rightarrow 0^+} \int\frac{\dd[4]{ p}}{\qty(2\pi)^42\omega_{\vb p}}\e^{\im \qty(x- y)\cdot p}\left(\frac{-2\omega_{\vb p}+2\im\epsilon}{\qty(p^0-\omega_{\vb p}+\im\epsilon)\qty(-p^0-\omega_{\vb p}+\im\epsilon)}\right)
\end{align}
to simplify further our result we need to say something of the $\epsilon\rightarrow0^+$ limit. It \textbf{does not} commute with the integral, why you may ask? The integrand has two poles, at $p^0 = \pm\qty(\omega_{\vb p}-\im\epsilon)$, 
if we take the $\epsilon$ limit before the integral then the integrand contains $\qty(p^0-\omega_{\vb p})^{-1}\qty(p^0+\omega_{\vb p})^{-1}$, as we're integrating $p^0$ over $\mathbb R$, the integrant is divergent, and the 
remaining integral does not converge or even is definable. But, there is something we can do, notice that the integral splits into two contributions,
\begin{align}
    \Delta_{\textnormal{F}}\qty(x,y) &=\frac{1}{\im}\lim\limits_{\epsilon\rightarrow 0^+} \int\frac{\dd[4]{ p}}{\qty(2\pi)^4}\frac{\e^{\im \qty(x- y)\cdot p}}{\qty(p^0-\omega_{\vb p}+\im\epsilon)\qty(-p^0-\omega_{\vb p}+\im\epsilon)}\nonumber\\
    &\quad\quad\quad-\lim\limits_{\epsilon\rightarrow 0^+}\epsilon \int\frac{\dd[4]{ p}}{\qty(2\pi)^4\omega_{\vb p}}\frac{\e^{\im \qty(x- y)\cdot p}}{\qty(p^0-\omega_{\vb p}+\im\epsilon)\qty(-p^0-\omega_{\vb p}+\im\epsilon)}
\end{align}
the second contribution is zero, to see this is enough to show that the limit of the integral doesn't have a contribution as $\frac1\epsilon$ or lower,
\begin{align}
    \lim\limits_{\epsilon\rightarrow 0^+}\int\frac{\dd[4]{ p}}{\qty(2\pi)^4\omega_{\vb p}}\frac{\e^{\im \qty(x- y)\cdot p}}{\qty(p^0-\omega_{\vb p}+\im\epsilon)\qty(-p^0-\omega_{\vb p}+\im\epsilon)}\sim\mathcal O\qty(1)
\end{align}
for this we'll need some machinery of residue theorem. Consider $x^0-y^0>0$, the opposite case is analogous, then the integral is convergent even for $\mathfrak{Im}\qty[p^0]<0$ due to the exponential suppression, even more, the contribution 
from considering the integral over a semi-circle going over only $\mathfrak{Im}\qty[p^0]\leq0$ of infinite radius is zero, due to the exponential suppression. That is, we can analytically continue the integral over a closed contour in the complex 
$p^0$ plane, starting at $p^0=-\infty$ going all the all to $+\infty$ then coming back to $-\infty$ through the lower half complex plane. Call this contour C,
\begin{align}
    \lim\limits_{\epsilon\rightarrow 0^+}\int\frac{\dd[4]{ p}}{\qty(2\pi)^4\omega_{\vb p}}\frac{\e^{\im \qty(x- y)\cdot p}}{\qty(p^0-\omega_{\vb p}+\im\epsilon)\qty(-p^0-\omega_{\vb p}+\im\epsilon)}=\lim\limits_{\epsilon\rightarrow 0^+}\int\frac{\dd[3]{\vb p}}{\qty(2\pi)^4\omega_{\vb p}}\oint\limits_{C}\dd{p^0}\frac{\e^{\im \qty(x- y)\cdot p}}{\qty(p^0-\omega_{\vb p}+\im\epsilon)\qty(-p^0-\omega_{\vb p}+\im\epsilon)}
\end{align}
as the integrand has only a pole in the lower half plane at $p^0=\omega_{\vb p}-\im\epsilon$, $C$ can be deformed to enclose this pole clockwise, thus, the integral is just minus the residue at this pole,
\begin{align}
    \lim\limits_{\epsilon\rightarrow 0^+}\int\frac{\dd[4]{ p}}{\qty(2\pi)^4\omega_{\vb p}}\frac{\e^{\im \qty(x- y)\cdot p}}{\qty(p^0-\omega_{\vb p}+\im\epsilon)\qty(-p^0-\omega_{\vb p}+\im\epsilon)}&=\lim\limits_{\epsilon\rightarrow 0^+}\int\frac{\dd[3]{\vb p}}{\qty(2\pi)^4\omega_{\vb p}}\frac{\e^{\im \qty(x- y)\cdot p}\qty(-2\pi\im)\qty(p^0-\omega_{\vb p}+\im\epsilon)}{\qty(p^0-\omega_{\vb p}+\im\epsilon)\qty(-p^0-\omega_{\vb p}+\im\epsilon)}\eval_{p^0=\omega_{\vb p}-\im\epsilon}\\
    \lim\limits_{\epsilon\rightarrow 0^+}\int\frac{\dd[4]{ p}}{\qty(2\pi)^4\omega_{\vb p}}\frac{\e^{\im \qty(x- y)\cdot p}}{\qty(p^0-\omega_{\vb p}+\im\epsilon)\qty(-p^0-\omega_{\vb p}+\im\epsilon)}&=\frac1\im\lim\limits_{\epsilon\rightarrow 0^+}\int\frac{\dd[3]{\vb p}}{\qty(2\pi)^3\omega_{\vb p}}\frac{\e^{\im \qty(x- y)\cdot p}}{\qty(-2\omega_{\vb p}+2\im\epsilon)}\eval_{p^0=\omega_{\vb p}-\im\epsilon}
\end{align}
here we can exchange the limit with the integral, as the integrand is well behaved in the integration domain even in the $\epsilon\rightarrow 0^+$ limit. Due to this, there isn't any ``negative powers'' of $\epsilon$ in the limit, hence, 
as this integral is multiplied by $\lim\limits_{\epsilon\rightarrow 0^+}\epsilon$, the whole term is zero. With this we can focus on the remaining term,
\begin{align}
    \Delta_{\textnormal{F}}\qty(x,y) &=\frac{1}{\im}\lim\limits_{\epsilon\rightarrow 0^+} \int\frac{\dd[4]{ p}}{\qty(2\pi)^4}\frac{\e^{\im \qty(x- y)\cdot p}}{\qty(p^0-\omega_{\vb p}+\im\epsilon)\qty(-p^0-\omega_{\vb p}+\im\epsilon)}\\
    \Delta_{\textnormal{F}}\qty(x,y) &=\frac{1}{\im}\lim\limits_{\epsilon\rightarrow 0^+} \int\frac{\dd[4]{ p}}{\qty(2\pi)^4}\frac{\e^{\im \qty(x- y)\cdot p}}{-\qty(p^0)^2+\omega_{\vb p}^2-2\im\epsilon\omega_{\vb p}-\epsilon^2}
\end{align}
by the same reasoning as before, the contribution $-\epsilon^2$ doesn't affect in nothing the limit, as it does not produce any additional pole, neither it cross a pole or branch-cut in the limit. The same can be said for the 
multiplying factor $2\omega_{\vb p}$ of $-\im\epsilon$. As $\omega_{\vb p}>0$ it does not impact the pole in the limit, only the sign in relevant, as it sets the way the pole is approached in the limit, thus, we can neglect the 
$-\epsilon^2$ term and change the modulus of $\epsilon$ to arrive at,
\begin{align}
    \Delta_{\textnormal{F}}\qty(x,y) &=\frac{1}{\im}\lim\limits_{\epsilon\rightarrow 0^+} \int\frac{\dd[4]{ p}}{\qty(2\pi)^4}\frac{\e^{\im \qty(x- y)\cdot p}}{p^2+m^2-\im\epsilon}
\end{align}